\section{案例演练：三类典型任务如何选目标与算法}

\subsection{案例 A：对抗型任务（接近二人零和）}
如果任务形态类似对抗博弈（例如红蓝攻防模拟），优先目标应是低 exploitability 而非短期平均收益。可采用：
\begin{itemize}
    \item 自博弈主线训练；
    \item 周期性 best-response probing；
    \item 以 exploitability 和 worst-case utility 作为主验收指标。
\end{itemize}

\begin{importantbox}{选择依据}
当环境允许对手针对性学习你的策略时，最坏情况风险比平均分更能预测线上生存能力。
\end{importantbox}

\subsection{案例 B：人类协作型任务（典型非二人零和）}
例如谈判助手、多方排期助手、协同决策助手。此类任务应把 Population Best Response 绑定到真实用户群体，并强制包含文化/习惯差异数据切片。

推荐流程：
\begin{enumerate}
    \item 先训练行为模仿模型，覆盖不同用户子群体；
    \item 再做策略优化，目标是群体总体效用加权；
    \item 最后做跨群体回归，防止某个子群体被系统性伤害。
\end{enumerate}

\begin{knowledgebox}{分群建模建议}
至少按地区、行业、任务偏好、风险容忍度做分桶，否则 ``平均最优'' 往往意味着 ``局部严重失配''。
\end{knowledgebox}

\subsection{案例 C：工具调用型任务（代码/数据分析 Agent）}
这类任务中多 Agent 收益主要来自路由与并行验证，不必一上来做复杂协商。实践上可以采用：
\begin{itemize}
    \item Planner-Agent 负责拆分子任务；
    \item Specialist-Agents 分别处理检索、代码、验证；
    \item Verifier-Agent 做结果一致性检查并回传置信度。
\end{itemize}

其收益机理更接近 ``专家系统''，而不是 ``多方博弈''，因此评测重点也应转为正确率-时延-成本三角。

\begin{warningbox}{错误迁移}
把谈判博弈中的策略算法原样搬到工具调用任务，常常只会增加 token 消耗，不会增加有效正确率。
\end{warningbox}

\subsection{统一决策表：先分类，再选目标}
\begin{table}[H]
\centering
\begin{tabular}{p{3.1cm}p{3.4cm}p{3.2cm}p{3.2cm}}
\toprule
\textbf{任务类型} & \textbf{优先目标} & \textbf{首选方法族} & \textbf{主评测指标} \\
\midrule
对抗型 & 低 exploitability & self-play + regret 系 & worst-case utility \\
人类协作型 & population utility & imitation + inference scaling + RL & 跨群体收益与满意度 \\
工具调用型 & 正确率/时延比 & routing + verifier 架构 & pass@k、时延、成本 \\
\bottomrule
\end{tabular}
\caption{从任务结构到算法选择的最短路径}
\end{table}

\subsection{本章小结}
同样叫 ``多 Agent''，三类任务的优化目标几乎不同。先分类再建模，能避免大量无效实验。

